\scanpage{71}
-\mspage{122}-
\margin{n° 371}
\uline{mineurs d’ordre} $p$ (\uline{pour} $p\leq\inf(m,n)$) \uline{les déterminants des matrices carrées d’ordre} $p$ \uline{obtenues à partir des sous-matrices à} $p$ \uline{lignes et} $p$ \uline{colonnes de la matrice $X$}.

Etant donnée une sous-matrice de $\underline X$ à $p$ lignes et $p$ colonnes, les diverses matrices carrées que l’on en déduit ne différant que par l’ordre des lignes ou des colonnes ; leurs déterminants sont donc les mêmes \uline{au signe près}. \uline{Dans toute la suite de ce chapitre}, nous nous bornerons au cas où $\underline X$ est une matrice dont l’ensemble des indices des lignes est l’intervalle $\lbrack\!\lbrack1,m\rbrack\!\rbrack$ de $N$ et l’ensemble des indices des colonnes l’intervalle $\lbrack\!\lbrack1,n\rbrack\!\rbrack$ de $N$. Soit alors $H$ (resp. $K$) une partie de $p$ éléments de l’ensemble $\lbrack\!\lbrack1,m\rbrack\!\rbrack$ (resp. $\lbrack\!\lbrack1,n\rbrack\!\rbrack$) et considérons la sous-matrice de $\underline X$ obtenue en supprimant dans $\underline X$ les lignes d’indice $i\in\complement H$ et les colonnes d’indice $j\in\complement K$ ; nous désignerons par $X_{H,K}$ le déterminant de la matrice carrée obtenue à partir de la sous-matrice précédente en \uline{rangeant par ordre croissant} les indices de ses lignes et de ses colonnes. Soit $(i_h)_{1\leq i\leq p}$ (resp. $(j_k)_{1\leq k\leq p}$) la suite strictement croissante des éléments de $H$ (resp. $K$) ; la définition précédente équivaut donc à la relation
\[
\begin{split}
(9)\quad &(e_{i_1}\xi_{i_1j_1}+e_{i_2}\xi_{i_2j_1}+\ldots+e_{i_p}\xi_{i_pj_1})\wedge\ldots\wedge(e_{i_1}\xi_{i_1j_p}+e_{i_2}\xi_{i_2j_p}+\ldots+e_{i_p}\xi_{i_pj_p})=\\
&=X_{H,K}.e_{i_1}\wedge e_{i_2}\wedge\ldots\wedge e_{i_p}.
\end{split}
\]
PROPOSITION 5.- \uline{Soit} $M$ \uline{un A-module ayant une base} $(e_i)_{1\leq i\leq n}$ \uline{de} $n$ \uline{éléments, et pour tout} $p\leq n$, \uline{soit} $(e_H)$ \uline{la base correspondante de} $\bigwedge^pM$, $H$ \uline{parcourant l’ensemble des parties de} $\lbrack\!\lbrack1,n\rbrack\!\rbrack$ \uline{à} $p$ \uline{éléments} (§ 5, n°\inserted{\unsure{3}}, prop.8). \uline{Soit} $(x_j)_{1\leq j\leq p}$ \uline{une suite de} $p$ \uline{éléments de} $M$, \uline{et soit $X$ la matrice à} $n$ \uline{lignes et} $p$ \uline{colonnes dont les composantes} $\xi_{ji}$ $(1\leq i\leq n)$ \uline{de} $x_j$ \uline{sur la base} $(e_i)$ \uline{forment la colonne d’indice} $j$. \uline{Alors la composante du} $p$-\uline{vecteur décomposable} $x_1\wedge x_2\wedge\ldots\wedge x_p$ \uline{sur} $e_H$ \uline{est le mineur} $X_{H,J}$ \uline{de $X$, où} $J=\lbrack\!\lbrack1,p\rbrack\!\rbrack$ (que l’on écrit aussi simplement $X_H$).

Cela résulte aussitôt de la formule (9), où on fait $j_k=k$ pour $1\leq k\leq p$ et où
\margin{70}
